\section{The ledger: every claim, its grade, and its evidence}\label{app:ledger}

\begin{center}
\small
\begin{tabular}{@{}p{0.30\textwidth}p{0.10\textwidth}p{0.50\textwidth}@{}}
\toprule
claim & grade & evidence \\
\midrule
Thm~\ref{thm:traceless} $\Tr\Fr=0$ & \T & \lean{GapsAndTotals.redistribution\_traceless} \\
Lem~\ref{lem:blocks} block identities & \T & \lean{redistribution\_block12}, \lean{redistribution\_block21}, \lean{block\_decomposition} \\
Thm~\ref{thm:compat} $\Fr=0\iff[\Omega,\PP]=0$ & \T & \lean{redistribution\_zero\_iff\_commute} \\
Thm~\ref{thm:norm} norm identity & \T & \lean{redistribution\_norm}, \lean{half\_commutator\_norm}, \lean{commutator\_skew} \\
Ex.~\ref{ex:nonnormal} non-normal counterexample & \T/\V & \lean{norm\_needs\_selfadjoint.lean} (must fail) \\
\midrule
Thm~\ref{thm:qgt} $Q$ Hermitian, $g$ PSD & \T & Cauchy--Schwarz on $Z_\mu$; \lean{redistribution\_norm} \\
Rem.~\ref{rem:qgt-correction} $Q^\circ=2g$ is real & \T/\V & \gate{check\_D2\_m7.py}: literal imaginary part $0$ \\
Thm~\ref{thm:curv} closedness, periods & \T/\Cc & Chern--Weil \\
Thm~\ref{thm:bound} curvature bound & \T/\V & \gate{check\_D3\_m1.py} \\
\midrule
Thm~\ref{thm:weyl} Weyl & \Cc & Weyl (1912) \\
Thm~\ref{thm:riesz}(a) finite pair, union gap & \Cc/\V & Davis--Kahan; \gate{g04\_projector\_constant} \\
Thm~\ref{thm:riesz}(b) infinitesimal, constant 1 & \T/\V & Davis--Kahan $+$ limit; \gate{g04} \\
Thm~\ref{thm:riesz}(c) angle, $\pi/2$ at the endpoint & \T/\Cc & $\arcsin$; Albeverio--Motovilov \\
Thm~\ref{thm:tan2} rank one, $\abs{\tan 2\Theta}\le2\delta/d$ & \T/\V & proof above; \gate{g06\_projector\_tan2} \\
Rem.~\ref{rem:tan} off-diagonal $\tan\Theta\le\delta/d$, rank $k$ & \Cc/\V & Motovilov--Selin; \gate{g05\_tan\_theta} \\
Thm~\ref{thm:rankone} rank one, no gain & \T/\V & \gate{g02\_gap\_hierarchy} (a) \\
Thm~\ref{thm:ep} exceptional point & \T/\V & \gate{g02\_gap\_hierarchy} (b) \\
Thm~\ref{thm:winding} winding and sign & \T/\V & \gate{g02\_gap\_hierarchy} (c); \gate{check\_D6\_m8.py} \\
\midrule
Lem~\ref{lem:inertia} inertia identity & \T & expansion \\
Thm~\ref{thm:perpairs} per-pair capacity & \T/\V & \gate{g01\_shape\_spectrum} (g) \\
Thm~\ref{thm:sumrule} sum rule & \T/\V & \gate{g01\_shape\_spectrum} (f) \\
\midrule
Thm~\ref{thm:shape-sum} $\sigma_1^2+\sigma_2^2=3$ & \T/\V & \gate{g01\_shape\_spectrum} (a) \\
Thm~\ref{thm:shape-det} determinant, maximum & \T/\V & \gate{g01\_shape\_spectrum} (b),(c),(h) \\
Thm~\ref{thm:cone} linear isotropic cone & \T/\V & \gate{g01\_shape\_spectrum} (d) \\
Thm~\ref{thm:shape-holonomy} holonomy $-1$ & \T/\V & \gate{g01\_shape\_spectrum} (e) \\
Thm~\ref{thm:optic} the optic points move & \T/\V & \gate{g03\_mass\_breaking} \\
\midrule
Thm~\ref{thm:snell} Snell for conformal metrics & \T & first variation \\
\midrule
physical readings (\S\ref{sec:scope}) & \Dd/\Hh & not used by any theorem \\
\bottomrule
\end{tabular}
\end{center}

\section{Reproduction}\label{app:repro}

\paragraph{Kernel.} The finite algebraic core is \lean{EmergentTransport/GapsAndTotals.lean}
(18 theorems) together with \lean{EmergentTransport/Basic.lean} (11 theorems), checked by
\begin{verbatim}
python3 scripts/verify.py --direct
# OK: 29 theorems, axioms used: ['Classical.choice', 'Quot.sound', 'propext'],
# 6 false controls rejected
\end{verbatim}
The verifier audits the axioms of every named theorem (only the three standard ones are allowed,
which excludes \lean{sorryAx}, project axioms and \lean{native\_decide}), requires six deliberately
false statements to \emph{fail} for a mathematical reason, and records the SHA-256 of every source
file. On a machine where the Mathlib checkout is writable, the same script runs under
\lean{lake env lean} without the flag.

\paragraph{Gates.} Every number in \S\S\ref{sec:redistribution}--\ref{sec:shape} is produced by a
gate that writes JSON:
\begin{verbatim}
python3 scripts/gates/g01_shape_spectrum.py
python3 scripts/gates/g02_gap_hierarchy.py
python3 scripts/gates/g03_mass_breaking.py
python3 scripts/gates/g04_projector_constant.py
python3 scripts/gates/g05_tan_theta.py
python3 scripts/gates/g06_projector_tan2.py
\end{verbatim}
The first has 20 assertions on the shape$\to$distance differential, the second 13 on rank one,
the exceptional point and the sign, the third 20 on the optic points under mass breaking, the fourth 10 on the projector constants and the failure of the constant-1 finite form, the fifth 4 on the off-diagonal $\tan\Theta$ bound, the sixth 12 on the rank-one $\tan 2\Theta$ theorem.
The paper's figures of merit (the closed form of $\det M^*M$, the cone slope, the holonomy, the
per-pair identity) are read from \gate{results/gates/*.json}, never typed by hand.

\paragraph{The swarm records.} The proofs of \S\ref{sec:hierarchy}.2--\S\ref{sec:hierarchy}.5 and of
Theorem~\ref{thm:perpairs} were produced and checked in eight independent directions whose
reports and scripts are in \gate{results/swarm/}, with \gate{README.md} recording what each one
corrected. Three claims in this paper were corrected during that process and the superseded
statements are kept in \gate{notes/02\_PLAN.md}, not deleted.

\section{The classical inputs, stated}\label{app:classical}

For completeness, the two results imported in \S\ref{sec:hierarchy}:

\emph{Weyl.} For self-adjoint $A,A'$ and eigenvalues ordered increasingly,
$\abs{\lambda_i(A')-\lambda_i(A)}\le\norm{A'-A}$. Proof: the min--max characterisation of the
$i$-th eigenvalue makes it $1$-Lipschitz in $A$ in the operator norm.

\emph{Davis--Kahan.} Let a spectral cluster $\shape$ of a self-adjoint $A$ satisfy
$\operatorname{dist}(\shape,\operatorname{spec}A\setminus\shape)\ge\Gap$, and let $\PP,\PP'$ be
the orthogonal projectors onto $\shape$ and onto the corresponding part of the spectrum of
$A'=A+E$. The sine theorem states $\norm{\PP-\PP'}\le\norm{E}/\Gap$, with the denominator the gap
in the \emph{union} of the two spectra --- which is at least $\Gap-\norm{E}$ when $\delta<\Gap/2$.
That is Theorem~\ref{thm:riesz}(a); the derivative form (b) follows by applying it to
$(A(t),A(t+h))$ and letting $h\to0$, and the angle form (c) by monotonicity of $\arcsin$.

\emph{On the constant $\pi/2$.} The estimates in this literature are stated for the
\emph{angle} $\arcsin\norm{\PP-\PP'}$ rather than for $\norm{\PP-\PP'}$ itself, and that is where
$\pi/2$ lives. Albeverio and Motovilov \Cc{} prove a bound of the form
$\arcsin\norm{\PP-\QQ}\le M(\norm{V}/\Gap)$ for $\norm{V}<\Gap/2$, and show that the two subspaces
are guaranteed to be in the acute-angle case only for $\norm{V}<c\,\Gap$ with
$c=0.454169\ldots<1/2$: that is, the angle reaches $\pi/2$ strictly before the stable range ends.
This is exactly the content of Theorem~\ref{thm:riesz}(c) expressed in the literature's
normalisation, and it is the source we attribute the angle form to. For \emph{off-diagonal}
perturbations the angle bound sharpens to Motovilov--Selin's $\tan\Theta\le\delta/d$, which never
reaches $\pi/2$; see Remark~\ref{rem:tan}, whose transcription of that bound is checked by gate
\gate{g05\_tan\_theta}.

We could \emph{not} verify that Kato's book states the constant $\pi/2$. The zbMATH review of the
1966 edition (Zbl 0148.12601) summarises Chapter II as finite-dimensional analytic perturbation
theory with perturbation series for eigenprojectors, a reduction method for semisimple
eigenvalues, and estimates of convergence radii --- without naming a constant --- and the sections
themselves were not available to us. The attribution to Kato is therefore dropped, not as a
disagreement but for want of evidence; the mathematics used in \S\ref{sec:hierarchy} is
Theorem~\ref{thm:riesz}, resting on Davis--Kahan for the finite-pair bound, on the limit for the
derivative, and on the elementary inequality $\arcsin x\le(\pi/2)x$ for the angle.

